\documentclass[draft]{beamer}

\usepackage[T1]{fontenc}
\usepackage[utf8]{inputenc}
\usepackage{textcomp}
\usepackage{color}
\usepackage{contour}
\usepackage{ulem}
\usepackage{natbib}
\usepackage{pgf}
\usepackage{xspace}
\xspaceaddexceptions{]\}}
\usepackage{multirow}
\usepackage{tikz}
\usetikzlibrary{calc,positioning,arrows}
\usepackage{ulem}
\usepackage{hyperref}
\usepackage[absolute,overlay]{textpos}
% \usepackage{beamerthemeshadow}
\usetheme{Warsaw}
\usepackage{lmodern}
\usepackage{mathtools}
% \everymath{\displaystyle}
\usepackage{algorithm}
\usepackage{algorithmic}
\usepackage{pbox}
% \usepackage[final]{listingsutf8}
\usepackage[french]{babel}
\usepackage[final]{listings}
\usepackage{courier}
\usepackage[outdir=./]{epstopdf}
\usepackage{ifdraft}
\setkeys{Gin}{draft=false}% force the display of graphics, even for draft mode

\lstset{
  language=Java,
  basicstyle=\lst@ifdisplaystyle\scriptsize\fi\ttfamily,
  commentstyle=\color[RGB]{63,127,95}\bfseries,
  keywordstyle=\color[RGB]{127,0,85},
  morekeywords={int,char,double,float,unsigned,void,bool},
  % Highlight some keywords
  classoffset=1,
  morekeywords={>,<},
  otherkeywords={>,<},
  keywordstyle=\color[RGB]{0,0,0}\bfseries,
  classoffset=0,
  % End highlight
  numbers=left,
  stepnumber=1,
  showstringspaces=false,
  tabsize=2,
  breaklines=false,
  escapeinside={§}{§},% similar to $, but not $ to allow math inside
  literate=%
    {à}{{\`a}}1
    {é}{{\'e}}1
    {è}{{\`e}}1
    {ê}{{\^e}}1
    {ï}{{\"i}}1
}
\definecolor{errorcolor}{RGB}{254,67,161}
\makeatletter
\def\uwave{\bgroup\markoverwith{\lower4\p@\hbox{\sixly \textcolor{errorcolor}{\char58}}}\ULon}
\font\sixly=lasy8
\makeatother
\newcommand{\javaerror}[1]{%
  \uwave{#1}%
}

\usetikzlibrary{positioning}
\makeatletter
\def\normaljustify{%
  \let\\\@centercr
  \rightskip\z@skip
  \leftskip\z@skip%
  \parfillskip=0pt plus 1.0fil\relax
}
\makeatother

% \join{-}{1,2,3} = 1-2-3
\makeatletter
\DeclareRobustCommand\join[2]{%
  \def\first{}%
  \def\list{#2}%
%   \@onelevel@sanitize\list%
  \@for\reserved:=\list\do{%
    \ifx\first\empty%
      \def\first{done}%
    \else%
      #1%
    \fi%
    \reserved%
  }
}
\makeatother

\mode<presentation>{
  \usetheme{Berlin}
  \setbeamercovered{invisible}
  \defbeamertemplate*{footline}{shadow theme}{
    \leavevmode
    \hbox{
      \begin{beamercolorbox}[wd=\paperwidth,ht=2.5ex,dp=1.125ex,leftskip=.3cm,rightskip=.3cm plus1fil]{title in head/foot}%
	\usebeamerfont{title in head/foot}\insertshorttitle%
	\hfill%
	\href{http://creativecommons.org/publicdomain/zero/1.0/}{\includegraphics[height=6pt]{../cc-zero}}%
	\qquad%
	\insertframenumber{}/{}\inserttotalframenumber%
      \end{beamercolorbox}
    }
    \hbox{
      \begin{beamercolorbox}[wd=\paperwidth,ht=2.5ex,dp=1.125ex,leftskip=.3cm plus1fil,rightskip=.3cm]{author in head/foot}
	\usebeamerfont{author in head/foot}\insertshortauthor\hfill\insertshortinstitute
      \end{beamercolorbox}
    }
    \vskip0pt
  }
%   \addtobeamertemplate{frametitle}{}{
%     \begin{textblock*}{30mm}(.8\textwidth,-7mm)
%       \includegraphics[height=6mm,keepaspectratio]{figures/logo-fbk}
%       \hfill
%       \includegraphics[height=6mm,keepaspectratio]{figures/logo-ict}
%     \end{textblock*}
%   }
  \setbeamerfont*{itemize/enumerate body}{size=\normalsize}
  \setbeamerfont*{itemize/enumerate subbody}{size=\footnotesize,shape=\itshape}
%   \setbeamertemplate{frametitle}[shadow]
}

\AtBeginSection[]{
  \begin{frame}
  \vfill
  \centering
  \begin{beamercolorbox}[sep=8pt,center,shadow=true,rounded=true]{title}
    \usebeamerfont{title}\insertsectionhead\par%
  \end{beamercolorbox}
  \vfill
  \end{frame}
}

\newcommand{\email}[1]{\href{mailto:#1}{\nolinkurl{#1}}}
\DeclareRobustCommand{\todo}[1]{\ifdraft{\textbf{\textcolor{red}{[#1]}}}{}}
\newcommand\credits[2][]{%
  \ifx&#1&%
  \def\x{}
  \else%
  \def\x{{#1}: }
  \fi%
  \begingroup%
  \renewcommand\thefootnote{}\footnote{{\x}{\textcopyright}#2}%
  \addtocounter{footnote}{-1}%
  \endgroup%
}
\def\newblock{\hskip .11em plus .33em minus .07em} %trick to avoid undefined command with bibtex
\def\randomSeed{1138}

\title{Java Generics\todo{40min}}
\author[Matthieu Vergne\qquad\email{matthieu.vergne@meritis.fr}]{
  \texorpdfstring{
    \textbf{Matthieu Vergne}
  }{Matthieu Vergne}
}
\institute[Meritis PACA]{
  Les Algorithmes - Aristote B\\
  2000 Route des Lucioles\\
  06901  Sophia Antipolis Cedex\\
  ~\\
  \includegraphics[height=15mm,keepaspectratio]{../logo-meritis.jpg}
}
\date{TBD}

\begin{document}

\begin{frame}
\titlepage
\end{frame}

\begin{frame}
\frametitle{Outline}
\tableofcontents[hideallsubsections]
\end{frame}

\section{Introduction}
\subsection{}

\begin{frame}
\frametitle{Résumé}
\begin{itemize}
 \item \todo{}
\end{itemize}
\end{frame}

\begin{frame}[fragile]
\frametitle{Contexte pre-Java 5}
\begin{onlyenv}<+>
\begin{lstlisting}
void doSomething() {
  List numbers = new List();
  numbers.add(123);
  numbers.add(456);
  String result = concat(numbers);
  System.out.println(result);
}

String concat(List strings) {
  String result = "";
  for(int i = 0; i < strings.size(); i++) {
    result += (String) strings.get(i);
  }
  return result;
}
\end{lstlisting}
\end{onlyenv}
\begin{onlyenv}<+>
\begin{lstlisting}
void doSomething() {
  List numbers = new List();
  numbers.add(123);
  numbers.add(456);
  String result = concat(numbers); // Compile
  System.out.println(result);
}

String concat(List strings) {
  String result = "";
  for(int i = 0; i < strings.size(); i++) {
    result += (String) strings.get(i); // Mais échoue ici
  }
  return result;
}
\end{lstlisting}
\end{onlyenv}
\end{frame}

\begin{frame}
\frametitle{Constat}
\begin{alertblock}{Problèmes}
 \begin{itemize}
  \item Contrôle à l'exécution
  \item Casts systématiques
  \item Classes spécifiques possibles, mais trop répétitives
 \end{itemize}
\end{alertblock}
\begin{exampleblock}{Solution: generics}
 Les generics Java sont un moyen de dire au compilateur le type attendu, de façon à ce qu'il effectue lui-même ces contrôles avant la moindre exécution :
 \begin{itemize}
  \item Contrôle à la compilation $\Rightarrow$ visible immédiatement
  \item Type connu $\Rightarrow$ cast inutile
  \item N'importe quel type $\Rightarrow$ classe spécifique inutile
 \end{itemize}
\end{exampleblock}
\end{frame}

\begin{frame}[fragile]
\frametitle{Example}
\begin{onlyenv}<+>
Avant :
\begin{lstlisting}
void doSomething() {
  List numbers = new List();
  numbers.add(123);
  numbers.add(456);
  String result = concat(numbers); // Compile
  System.out.println(result);
}

String concat(List strings) {
  String result = "";
  for(int i = 0; i < strings.size(); i++) {
    result += (String) strings.get(i); // Mais échoue ici
  }
  return result;
}
\end{lstlisting}
\end{onlyenv}
\begin{onlyenv}<+>
Après :
\begin{lstlisting}
void doSomething() {
  List<Integer> numbers = new List<>(); // Fournit Integer
  numbers.add(123);
  numbers.add(456);
  String result = §\javaerror{concat(numbers)}§; // Ne compile pas
  System.out.println(result);
}

String concat(List<String> strings) { // Demande String
  String result = "";
  for(int i = 0; i < strings.size(); i++) {
    result += strings.get(i); // Plus besoin de cast
  }
  return result;
}
\end{lstlisting}
\end{onlyenv}
\end{frame}

\section{Cas d'usage : processing distribué}

\subsection{Alignement des champs d'une classe}

\begin{frame}[fragile]
\frametitle{Contexte}
\begin{block}{Création d'une paire d'objets de même type}
 Deux éléments de même type doivent être renvoyés par une méthode.
 Java n'a pas de \lstinline{return} multiple, on doit donc utiliser un conteneur.
 Une \lstinline{List} n'a pas de taille fixe, on fera donc une classe \lstinline{Pair}.
 Il faut pouvoir traiter différents types pour l'utiliser dans différents contextes.
\end{block}
\begin{lstlisting}
class Pair {
  Object first;
  Object second;
}
\end{lstlisting}
\end{frame}

\begin{frame}[fragile]
\frametitle{Problème}
\begin{alertblock}{Manque de contrôles}
 \begin{itemize}
  \item Pas de contrainte de type entre les champs
  \item Cast obligatoire à la lecture
 \end{itemize}
\end{alertblock}
\begin{lstlisting}
// int + String = non voulu mais compile
Pair pair = new Pair();
pair.first = 123;
pair.second = "bar";

// Lecture nécessite cast
Pair strings = new Pair();
strings.first = "foo";
String first = (String) strings.first;
\end{lstlisting}
\end{frame}

\begin{frame}[fragile]
\frametitle{Solution}
\begin{exampleblock}{Generics pour les champs}
\begin{lstlisting}[numbers=none]
class Pair<T> {
  T first;
  T second;
}
\end{lstlisting}
\end{exampleblock}
\begin{lstlisting}
Pair<Integer> pair = new Pair<>(); // Demande Integer
pair.first = 123;
pair.second = §\javaerror{"bar"}§; // Ne compile pas

Pair<String> strings = new Pair<>();
strings.first = "foo";
String first = strings.first; // Pas de cast
\end{lstlisting}
\end{frame}

\subsection{Alignement des méthodes d'une interface}

\begin{frame}[fragile]
\frametitle{Contexte}
\begin{block}{Permettre différentes implémentations de \lstinline{Pair}}
 On souhaite lire les valeurs de différentes manières : valeurs fixes, variables, reprises d'autres objets, etc.
 On définit l'interface commune, chaque implémentation décidera de quoi mettre derrière.
\end{block}
\begin{lstlisting}
interface Pair {
  Object getFirst();
  Object getSecond();
}
\end{lstlisting}
\end{frame}

\begin{frame}[fragile]
\frametitle{Problème}
\begin{alertblock}{Manque de contrôles}
 \begin{itemize}
  \item Pas de contrainte de type entre les méthodes
  \item Cast obligatoire à l'appel
 \end{itemize}
\end{alertblock}
\begin{lstlisting}
// int + String = non voulu mais compile
Pair pair = new Pair() {
  public Object getFirst() {return 123;}
  public Object getSecond() {return "bar";}
};

// Lecture nécessite cast
Pair strings = new Pair() {...};
String first = (String) strings.getFirst();
String second = (String) strings.getSecond();
\end{lstlisting}
\end{frame}

\begin{frame}[fragile]
\frametitle{Solution}
\begin{onlyenv}<+>
\begin{exampleblock}{Generics pour les méthodes}
\begin{lstlisting}[numbers=none]
interface Pair<T> {
  T getFirst();
  T getSecond();
}
\end{lstlisting}
\end{exampleblock}
\begin{lstlisting}
Pair<Integer> pair = new Pair<Integer>() { // Demande Integer
  public Integer getFirst() {return 123;}
  public Integer getSecond() {return §\javaerror{"bar"}§;} // Ne compile pas
};

Pair<String> strings = new Pair<String>() {...};
String first = strings.getFirst();   // Pas de cast
String second = strings.getSecond(); // Pas de cast
\end{lstlisting}
\end{onlyenv}
\begin{onlyenv}<+>
On pourra utiliser l'implémentation suivante :
\begin{lstlisting}
class FinalPair<T> implements Pair<T> {
  private final T first;
  private final T second;

  public FinalPair(T first, T second) {
    this.first = first;
    this.second = second;
  }

  public T getFirst() {return first;}
  public T getSecond() {return second;}
}
\end{lstlisting}
\end{onlyenv}
\end{frame}

\subsection{Consommation de generics}

\begin{frame}[fragile]
\frametitle{Contexte}
\begin{block}{Affichage d'une \lstinline{Pair}}
 On veut un \lstinline{Pair.toString()} qui affiche ses 2 éléments.
 \lstinline{Pair} étant une interface, on ne peut directement l'y implémenter.
 On y crée une méthode \lstinline{static} utilisable dans n'importe quelle implémentation.
\end{block}
\begin{lstlisting}
static String toString(Pair<Object> p) {
  return String.format("(%s,%s)", p.getFirst(), p.getSecond());
}
\end{lstlisting}
\end{frame}

\begin{frame}[fragile]
\frametitle{Problème}
\begin{alertblock}{N'accepte que des \lstinline{Pair<Object>}}
 L'héritage ne s'applique pas automatiquement dans les generics :
 \begin{itemize}
  \item \lstinline{<Object>} n'accepte que \lstinline{<Object>}
  \item \lstinline{<String>} n'accepte que \lstinline{<String>}
  \item etc.
 \end{itemize}
\end{alertblock}
\begin{lstlisting}
// Aucun type ne compile en dehors de Object
§\javaerror{toString}§(new FinalPair<Integer>(1, 2));
§\javaerror{toString}§(new FinalPair<String>("a", "b"));
\end{lstlisting}
\end{frame}

\begin{frame}[fragile]
\frametitle{Solution}
\begin{onlyenv}<+>
\begin{exampleblock}{Wildcard \lstinline{\?}}
\begin{lstlisting}[numbers=none]
static String toString(Pair<?> p) {
  return String.format("(%s,%s)", p.getFirst(), p.getSecond());
}
\end{lstlisting}
\end{exampleblock}
\begin{lstlisting}
toString(new FinalPair<Integer>(1, 2));
toString(new FinalPair<String>("a", "b"));
\end{lstlisting}
\texttt{\bf\color{red}/!\textbackslash} \lstinline{Pair p} (\textit{raw type}) revient à désactiver tout contrôle (rétrocompatibilité pré-Java 5).
Ne jamais utiliser !
\end{onlyenv}
\begin{onlyenv}<+>
Exemple:\todo{Préciser les raisons exactes}
\begin{lstlisting}
List<HashMap> list = ...
List<Map<?, ?>> linkedList = new LinkedList<>(list); // OK
\end{lstlisting}
\begin{lstlisting}
List<HashMap> list = ...
List<Map<?, ?>> linkedList = new LinkedList<Map<?, ?>>(list); // KO
\end{lstlisting}
Due au ``raw type" (HashMap).
Si on corrige (\lstinline{HashMap< ?, ?>}) pas de soucis, mais le \lstinline{…} doit effectivement renvoyer quelque chose de compatible.
Un parseur par exemple, qui prendrait un \lstinline{HashMap.class} pour instancier, ne pourrait pas.
\todo{Parler de la traduction, dans notre cas un \lstinline{list.stream.map(...).collect(...)} permet de typer correctement}
\end{onlyenv}
\begin{onlyenv}<+>
On pourra compléter notre implémentation comme suit :
\begin{lstlisting}
class FinalPair<T> implements Pair<T> {
  ...
  public String toString() {return Pair.toString(this);}
}
\end{lstlisting}
\end{onlyenv}
\end{frame}

\subsection{Alignement des types d'une méthode}

\begin{frame}[fragile]
\frametitle{Contexte}
\begin{block}{Abstraction des opérations sur une paire d'éléments}
 On transforme 2 éléments en 1 seul, comme sommer des entiers, concaténer des \lstinline{Strings}, etc.
 Les 3 éléments ont un type quelconque, mais le même type pour permettre des exécutions à la chaîne.
 On pourra par exemple sommer toute une liste d'entiers en les sommant deux à la fois.
\end{block}
\begin{lstlisting}
interface Computer {
  Object compute(Object o1, Object o2);
}
\end{lstlisting}
\end{frame}

\begin{frame}[fragile]
\frametitle{Problème}
\begin{alertblock}{Manque de contrôles}
 \begin{itemize}
  \item Pas de contrainte de type entre les éléments
  \item Cast obligatoire dans l'implémentation
  \item Cast obligatoire à l'appel
 \end{itemize}
\end{alertblock}
\begin{lstlisting}
Computer ints = new Computer() {
  public Object compute(Object o1, Object o2) {
    return ((int) o1) + ((int) o2); // Casts
  }
};
int i = (int) ints.compute(1, 2); // Cast

// Mauvais types partout, mais compile
String s = (String) ints.compute("a", new Object());
\end{lstlisting}
\end{frame}

\begin{frame}[fragile]
\frametitle{Solution}
\begin{exampleblock}{Generics pour tous les éléments}
\begin{lstlisting}[numbers=none]
interface Computer<T> {
  T compute(T t1, T t2);
}
\end{lstlisting}
\end{exampleblock}
\begin{lstlisting}
Computer<Integer> ints = new Computer<Integer>() {
  public Integer compute(Integer i1, Integer i2) {
    return i1 + i2; // Sans cast
  }
};
int i = ints.compute(1, 2); // Sans cast

int j = ints.§\javaerror{compute}§("a", 2);  // Mauvaise entrée
String l = §\javaerror{ints.compute(1, 2)}§; // Mauvaise sortie
\end{lstlisting}
\end{frame}

\subsection{Alignement avec des generics d'autres classes}

\begin{frame}[fragile]
\frametitle{Contexte}
\begin{block}{On s'appuie sur l'interface \lstinline{Pair<T>}}
 \lstinline{Pair<T>} met déjà en place du contrôle à la compilation.
 On remplace donc les 2 entrées de \lstinline{compute} par une unique \lstinline{Pair<T>}.
\end{block}
\begin{lstlisting}
interface Computer {
  Object compute(Pair<?> pair);
}
\end{lstlisting}
\end{frame}

\begin{frame}[fragile]
\frametitle{Problème}
\begin{alertblock}{Contrôle incomplet}
 L'usage de \lstinline{Pair<T>} garantie l'alignement des entrées.
 Mais on n'en connait pas le type, obligeant le cast dans l'implémentation.
 La sortie n'est pas non plus connue, obligeant le cast aussi.
\end{alertblock}
\begin{lstlisting}
Computer ints = new Computer() {
  public Object compute(Pair<?> pair) {
    return (int) pair.getFirst()   // Cast entrée
         + (int) pair.getSecond(); // Cast entrée
  }
};
int i = (int) ints.compute(new FinalPair<>(1, 2)); // Cast sortie

// Mauvais type de sortie = compile quand même
String s = (String) ints.compute(new FinalPair<>(1, 2));
\end{lstlisting}
\end{frame}

\begin{frame}[fragile]
\frametitle{Solution}
\begin{onlyenv}<+>
\begin{exampleblock}{Generics pour entrée et sortie}
\begin{lstlisting}[numbers=none]
interface Computer<T> {
  T compute(Pair<T> pair);
}
\end{lstlisting}
\end{exampleblock}
\begin{lstlisting}
Computer<Integer> ints = new Computer<Integer>() {
  public Integer compute(Pair<Integer> pair) {
    return pair.getFirst() + pair.getSecond(); // Sans cast
  }
};
int i = ints.compute(new FinalPair<>(1, 2)); // Sans cast

// Mauvais type de sortie = ne compile pas
String l = §\javaerror{ints.compute(new FinalPair<>(1, 2))}§;
\end{lstlisting}
\end{onlyenv}
\begin{onlyenv}<+>
On pourra lui donner des implémentations spécifiques :
\begin{lstlisting}
// Somme d'entiers
class SumComputer implements Computer<Integer> {
  public Integer compute(Pair<Integer> pair) {
    return pair.getFirst() + pair.getSecond();
  }
}
new SumComputer().compute(new FinalPair<Integer>(1, 2));
new SumComputer().§\javaerror{compute}§(new FinalPair<String>("a", "b"));
// Concatènation de Strings
class ConcatComputer implements Computer<String> {
  public String compute(Pair<String> pair) {
    return pair.getFirst() + pair.getSecond();
  }
}
new ConcatComputer().§\javaerror{compute}§(new FinalPair<Integer>(1, 2));
new ConcatComputer().compute(new FinalPair<String>("a", "b"));
\end{lstlisting}
\end{onlyenv}
\begin{onlyenv}<+>
On pourra lui donner des implémentations génériques :
\begin{lstlisting}
// Extraction du premier élément d'une Pair quelconque
class ExtractComputer<T> implements Computer<T> {
  public T compute(Pair<T> pair) {
    return pair.getFirst();
  }
}
\end{lstlisting}
Le point clé est l'ajout d'un generics pour assurer le lien des types d'entrée-sortie sans lui donner de valeur particulière.
Cela oblige néanmoins à créer une instance pour chaque type de \lstinline{Pair} :
\begin{lstlisting}
new ExtractComputer<Integer>().compute(new FinalPair<>(1, 2));
new ExtractComputer<String>().compute(new FinalPair<>("a", "b"));
\end{lstlisting}
\end{onlyenv}
\end{frame}

\subsection{Alignement de scope méthode}

\begin{frame}[fragile]
\frametitle{Contexte}
\begin{block}{Opérations distribuées}
 Les opérations ne doivent pas être exécutées localement.
 Elles doivent être exécutées sur plusieurs machines cibles.
 Chaque machine doit pouvoir traiter n'importe quelle demande.
 Contraindre une cible à 1 seul type (\lstinline{Target<T>}) n'est donc pas souhaitable.
\end{block}
\begin{lstlisting}
interface Target {
  void dispatch(Pair<?> pair, Computer<?> computer);
}
\end{lstlisting}
\end{frame}

\begin{frame}[fragile]
\frametitle{Problème}
\begin{alertblock}{Inexploitable}
 La \lstinline{Pair} et le \lstinline{Computer} son typés mais inconnus (\lstinline{<?>}).
 Impossible de les faire interagir, car ils pourraient être incompatibles (e.g. \lstinline{Pair<Integer>} et \lstinline{Computer<String>}).
\end{alertblock}
\begin{lstlisting}
Target target = new Target() {
  public void dispatch(Pair<?> pair, Computer<?> computer) {
    // Cette ligne devrait être envoyée à la cible.
    // On l'exécute ici par simplicité.
    computer.§\javaerror{compute}§(pair); // Ne compile pas
  }
};
\end{lstlisting}
\end{frame}

\begin{frame}[fragile]
\frametitle{Solution}
\begin{onlyenv}<+>
\begin{exampleblock}{Generics commun sur la méthode}
\begin{lstlisting}[numbers=none]
interface Target {
  <T> void dispatch(Pair<T> pair, Computer<T> computer);
}
\end{lstlisting}
\end{exampleblock}
\begin{lstlisting}
Target target = new Target() {
  public <T> void dispatch(Pair<T> pair, Computer<T> computer) {
    // Cette ligne devrait être envoyée à la cible.
    // On l'exécute ici par simplicité.
    computer.compute(pair);
  }
};
\end{lstlisting}
\end{onlyenv}
\begin{onlyenv}<+>
On s'assurera que la distribution est correctement contrainte :
\begin{lstlisting}
FinalPair<Integer> intPair = new FinalPair<>(1, 2);
target.dispatch(intPair, new SumComputer());
target.§\javaerror{dispatch}§(intPair, new ConcatComputer()); // Integer != String
target.dispach(intPair, new ExtractComputer<>());

FinalPair<String> strPair = new FinalPair<>("a", "b");
target.§\javaerror{dispatch}§(strPair, new SumComputer());    // String != Integer
target.dispatch(strPair, new ConcatComputer());
target.dispatch(strPair, new ExtractComputer<>());
\end{lstlisting}
\end{onlyenv}
\begin{onlyenv}<+>
S'applique aussi aux constructeurs :
\begin{lstlisting}
class Dispatcher {
  Runnable preparedDispatch;
  
  <T> Dispatcher(Pair<T> pair, Computer<T> computer, Target target) {
    preparedDispatch = () -> target.dispatch(pair, computer);
  }
  
  void run() {
    preparedDispatch.run();
  }
}
\end{lstlisting}
Generics inutile au niveau de la classe, car aligne uniquement des paramètres du constructeur.
\end{onlyenv}
\end{frame}

\subsection{Défaut}

\begin{frame}[fragile]
\frametitle{Contexte}
\todo{? extends X}
\todo{? super X}
\todo{cheat sheet, diagramme question-réponse}
\begin{block}{\todo{todo}}
 \todo{todo}
\end{block}
\begin{lstlisting}
// TODO
\end{lstlisting}
\end{frame}

\begin{frame}[fragile]
\frametitle{Problème}
\begin{alertblock}{\todo{todo}}
 \todo{todo}
\end{alertblock}
\begin{lstlisting}
// TODO
\end{lstlisting}
\end{frame}

\begin{frame}[fragile]
\frametitle{Solution}
\begin{exampleblock}{\todo{todo}}
\begin{lstlisting}[numbers=none]
// TODO
\end{lstlisting}
\end{exampleblock}
\begin{lstlisting}
// TODO
\end{lstlisting}
\end{frame}

\section{Usages avancés}

\subsection{Cas 1}

\begin{frame}[fragile]
\frametitle{Contexte}
\begin{block}{\todo{todo}}
 \todo{todo}
\end{block}
\begin{lstlisting}
// TODO
\end{lstlisting}
\end{frame}

\begin{frame}[fragile]
\frametitle{Problème}
\begin{alertblock}{\todo{todo}}
 \todo{todo}
\end{alertblock}
\begin{lstlisting}
// TODO
\end{lstlisting}
\end{frame}

\begin{frame}[fragile]
\frametitle{Solution}
\begin{exampleblock}{\todo{todo}}
\begin{lstlisting}[numbers=none]
// TODO
\end{lstlisting}
\end{exampleblock}
\begin{lstlisting}
// TODO
\end{lstlisting}
\end{frame}

\subsection{Défaut}

\begin{frame}[fragile]
\frametitle{Contexte}
\todo{\lstinline{Pair<Collection>}, mais \lstinline{List getFirst()}}
\todo{et-commercial pour combinaison d'interfaces + abstract}
\begin{block}{\todo{todo}}
 \todo{todo}
\end{block}
\begin{lstlisting}
// TODO
\end{lstlisting}
\end{frame}

\begin{frame}[fragile]
\frametitle{Problème}
\begin{alertblock}{\todo{todo}}
 \todo{todo}
\end{alertblock}
\begin{lstlisting}
// TODO
\end{lstlisting}
\end{frame}

\begin{frame}[fragile]
\frametitle{Solution}
\begin{exampleblock}{\todo{todo}}
\begin{lstlisting}[numbers=none]
// TODO
\end{lstlisting}
\end{exampleblock}
\begin{lstlisting}
// TODO
\end{lstlisting}
\end{frame}

\section{Pièges}

\subsection{Cas précédents}

\begin{frame}
\frametitle{Résumé}
\begin{table}
  \begin{tabular}{|c|c|c|l}
    \cline{1-3}
    Scope                    & Cas             & Vu \\
    \cline{1-3}
    \multirow{2}{*}{Méthode} & Param-param     & x  & $\leftarrow$ Piège\\
    \cline{2-3}
    ~                        & Param-return    & o  \\
    \cline{1-3}
    N méthodes               & Méthode-méthode & o  \\
    \cline{1-3}
    \multirow{2}{*}{Classe}  & 1 méthode       & o  \\
    \cline{2-3}
    ~                        & N méthodes      & o  \\
    \cline{1-3}
  \end{tabular}
\end{table}
\end{frame}

\subsection{Param-param}

\begin{frame}[fragile]
\frametitle{Idée}
\begin{block}{Comparateur de fichier intelligent (diff)}
 Une application graphique capable de comparer différents types de fichiers.
 L'interface affiche un fichier à gauche et l'autre, du même type, à droite.
 Elle met en avant les différences en fonction d'algorithmes de comparaison de votre cru.
\end{block}
\begin{lstlisting}
interface Displayer {
  void load(File file);
}

class TxtDisplayer implements Displayer {...}
class XmlDisplayer implements Displayer {...}
...

<D extends Displayer> // Generics pour forcer le même type
void feed(D leftDisplayer, D rightDisplayer);
\end{lstlisting}
\end{frame}

\begin{frame}[fragile]
\frametitle{Piège}
\begin{alertblock}{Le generics n'interdit pas les types parents}
 La résolution du type se fait au plus spécifique que possible.
 Cependant, elle abstrait tant que nécessaire.
 Au pire, elle prend au plus générique, même \lstinline{Object}.
\end{alertblock}
\begin{lstlisting}
// D = TxtDisplayer
feed(txtDisplayerLeft, txtDisplayerRight);

// D = XmlDisplayer
feed(xmlDisplayerLeft, xmlDisplayerRight);

// D = Displayer <- permet de compiler
feed(txtDisplayerLeft, xmlDisplayerRight);
\end{lstlisting}
\end{frame}

\begin{frame}[fragile]
\frametitle{Solution}
\begin{exampleblock}{Generics explicite}
 Il faut forcer le generics si on sait quel type on veut.
\end{exampleblock}
\begin{lstlisting}
// D = XmlDisplayer <- ne compile pas avec TxtDisplayer
this.<XmlDisplayer>§\javaerror{feed}§(txtDisplayerLeft, xmlDisplayerRight);
\end{lstlisting}
\begin{exampleblock}{Délégation à l'appelant}
 Si on ne connaît pas le type ici, c'est une méthode appelante qui doit le connaître.
\end{exampleblock}
\begin{lstlisting}
<D extends Displayer> void buildScreen(D left, D right) {
  feed(left, right);
}
this.<XmlDisplayer>§\javaerror{buildScreen}§(txtDisplayerLeft, xmlDisplayerRight);
\end{lstlisting}
\end{frame}

\begin{frame}[fragile]
\frametitle{Generics inutile ?}
Au même niveau : generics explicite + arguments bien choisis
\begin{lstlisting}
this.<XmlDisplayer>feed(xmlDisplayerLeft, xmlDisplayerRight);
\end{lstlisting}
Simplification possible : retrait du generics
\begin{lstlisting}
void feed(Displayer leftDisplayer, Displayer rightDisplayer) {...}
feed(xmlDisplayerLeft, xmlDisplayerRight);
\end{lstlisting}
Remplacer un argument par un \lstinline{TxtDisplayer} ne compile pas dans le premier cas mais compile dans le second.
L'ajout du generics permet de découpler l'information de type (generics) et de valeur (arguments).
Il faut cependant penser à l'expliciter pour que ce soit utile.
\end{frame}

\begin{frame}[fragile]
\frametitle{Piège valable sur Generics aussi}
\todo{Même section ou nouvelle ?}
\begin{lstlisting}
toString(Pair<Object> p) {...}

// Compile
toString(new FinalPair<>(1, 2));
toString(new FinalPair<>("a", "b"));
\end{lstlisting}
\begin{alertblock}{\lstinline{<>} prend ce qui l'arrange pour que ça compile}
 \begin{itemize}
  \item \lstinline{FinalPair<>(1, 2)} $\rightarrow$ \lstinline{<Integer>} ou parent
  \item \lstinline{FinalPair<>("a", "b")} $\rightarrow$ \lstinline{<String>} ou parent
  \item \lstinline{toString()} $\rightarrow$ uniquement \lstinline{<Object>}
  \item \lstinline{<Object>} est retenu pour les deux car seul compilable
 \end{itemize}
\end{alertblock}
\end{frame}

\begin{frame}[fragile]
\frametitle{Solution}
\begin{exampleblock}{Expliciter au moins une fois}
 Ne pas utiliser \lstinline{<>} si le type n'est donné nulle part dans l'instruction.
\end{exampleblock}
\begin{lstlisting}
// Si seul présent, <> devrait être explicité
§\javaerror{toString}§(new FinalPair<Integer>(1, 2));

// Précisé par la variable, donc <> acceptable
Pair<Integer> pair = new FinalPair<>(1, 2)
§\javaerror{toString}§(pair);

// Précisé par 1 argument, le second doit s'aligner
<T> toString(Pair<T> p1, Pair<T> p2) {...}
§\javaerror{toString}§(new FinalPair<Integer>(1, 2), §\javaerror{new FinalPair<>("a", "b")}§);
\end{lstlisting}
\end{frame}

\subsection{Défaut}

\begin{frame}[fragile]
\frametitle{Idée}
\todo{generics perdu à l'exécution : tout comme sans generics, ne permet pas plusieurs implémentations}
\todo{\lstinline{compute(Pair pair)} et \lstinline{compute(Pair<Object> pair)} acceptent \lstinline{new FinalPair<>(1,2)} alors que c'est un \lstinline{Pair<Integer>} : faux, la méthode impose \lstinline{Pair<Object>} , donc il interpréte \lstinline{new FinalPair<Object>(1,2)}, il l'accepte donc parce que \lstinline{Object} peut prendre n'importe quoi, y compris des entiers. Si on écrit explicitement \lstinline{new FinalPair<Integer>(1,2)} ça ne compile plus.}
\todo{dans lambda, scope méthode peut devoir être explicité, y compris avec this}
\todo{"interface Computer<T> {<U extends T> U compute(Pair<U> pair);}" et autres définitions pas toujours possibles à implémenter.}
\begin{block}{\todo{todo}}
 \todo{todo}
\end{block}
\begin{lstlisting}
// TODO
\end{lstlisting}
\end{frame}

\begin{frame}[fragile]
\frametitle{Piège}
\begin{alertblock}{\todo{todo}}
 \todo{todo}
\end{alertblock}
\begin{lstlisting}
// TODO
\end{lstlisting}
\begin{center}
 \todo{conclusion}
\end{center}
\end{frame}

\section{Conclusion}
\subsection{}

\begin{frame}
\frametitle{Résumé}
\todo{}
\end{frame}

\begin{frame}
\frametitle{Aller plus loin}
\begin{itemize}
 \item Generics (Oracle) : \url{https://docs.oracle.com/javase/tutorial/java/generics/index.html}
 \item \todo{}
\end{itemize}
\end{frame}

\begin{frame}
\begin{center}
Thanks for your attention.\\
~\\
Questions?\\
% \nocite{*}
\end{center}
\end{frame}

%%%%%%%%%%%%%%%%%%%%%%%%%%%%%%%%%%%%%%%%%%%%%%%%%%%%%%%%%%%%

\begin{frame}[fragile]
\frametitle{Scope Méthode}
\todo{Faire 1 slide/cas : contexte, problème, solution avec generics}
\begin{onlyenv}<+>
Même type de paramètres :
\begin{lstlisting}
<T> void f(T t1, T t2);

f(string1, string2); // OK
f(int1, int2);       // OK
f(string1, int2);    // Ne compile pas
\end{lstlisting}
\end{onlyenv}
\begin{onlyenv}<+>
Même type de sortie que d'entrée :
\begin{lstlisting}
<T> T f(T t);

String s2 = f(s1); // OK
int i2 = f(i1);    // OK
int i2 = f(s1);    // Ne compile pas
\end{lstlisting}
\end{onlyenv}
\begin{onlyenv}<+>
N'importe quelle combinaison :
\begin{lstlisting}
<T> T f(String s, T t1, int i, T t2);

String s3 = f("", s1, 0, s2); // OK
int i3 = f("", i1, 0, i2);    // OK
int i3 = f("", s1, 0, i2);    // Ne compile pas
\end{lstlisting}
\end{onlyenv}
\begin{onlyenv}<+>
Types spécifiques : \lstinline{extends}
\begin{lstlisting}
class A {...}
class B extends A {...}
class C extends B {...}

<T extends B> T f(T t1, T t2);

A a = f(new A(), new A()) // Ne compile pas, A pas un B
B b = f(new B(), new B()) // OK, T = B
C c = f(new C(), new C()) // OK, T = C qui est un B
B b = f(new B(), new C()) // OK, T = B pour accepter B et C
\end{lstlisting}
Si on enlevait T pour avoir \lstinline{B f(B b1, B b2)} directement, on retournerait toujours \lstinline{B}, même à la ligne 9.
\end{onlyenv}
\begin{onlyenv}<+>
Types spécifiques : \lstinline{super}
\todo{todo}
\end{onlyenv}
\begin{onlyenv}<+>
Plusieurs types :
\begin{lstlisting}
<T, U> U f(T t1, T t2, U u);

String s2 = f(i1, i2, s1); // OK
String i2 = f(s1, s2, i1); // OK
String s4 = f(s1, s2, s3); // OK, même si T = U
\end{lstlisting}
\end{onlyenv}
\begin{onlyenv}<+>
Points d'attention :
\begin{itemize}
 \item Generics pour 1 élément -> remplacer par Object
\end{itemize}
\end{onlyenv}
\end{frame}

\begin{frame}[fragile]
\frametitle{Méthodes multiples ?}
\begin{onlyenv}<+>
\begin{lstlisting}
class C {
 <T> void f1(T t); // Son propre type T
 <T> void f2(T t); // Son propre type T
}

C c = new C();
c.f1(string1); c.f2(string2); // OK
c.f1(int1); c.f2(int2);       // OK
c.f1(string1); c.f2(int2);    // OK aussi !
\end{lstlisting}
Les types ne sont pas liés entre méthodes.
Impossible de les lier sans savoir quand un changement de type est acceptable.
\end{onlyenv}
\begin{onlyenv}<+>
Ajouter un patron de méthode\footnote<.>{\url{https://fr.wikipedia.org/wiki/Patron_de_m\%C3\%A9thode}} pour lier les types :
\begin{lstlisting}
class C {
 <T> void f1(T t);               // Son propre type T
 <T> void f2(T t);               // Son propre type T
 <T> void f1_et_f2(T t1, T t2) { // Son propre type T
  f1(t1); // Utilise T de f1_et_f2 sur f1
  f2(t2); // Utilise T de f1_et_f2 sur f2
 }
}

C c = new C();
c.f1_et_f2(string1, string2); // OK
c.f1_et_f2(int1, int2);       // OK
c.f1_et_f2(string1, int2);    // Ne compile pas
\end{lstlisting}
\end{onlyenv}
\end{frame}

\begin{frame}[fragile]
\frametitle{Scope Classe}
\begin{onlyenv}<+>
Un type par instance :
\begin{lstlisting}
class C<T> {
 void f1(T t); // Type T commun
 void f2(T t); // Type T commun
}

C<String> c = new C<>();
c.f1(myString);c.f2(myString); // OK
c.f1(myInt);c.f2(myInt);       // Ne compilent pas

C<Integer> c = new C<>();
c.f1(myString);c.f2(myString); // Ne compilent pas
c.f1(myInt);c.f2(myInt);       // OK
\end{lstlisting}
\end{onlyenv}
\begin{onlyenv}<+>
N'importe quelle combinaison :
\begin{lstlisting}
<T> T f(String s, T t1, int i, T t2);

String s3 = f("", s1, 0, s2); // OK
int i3 = f("", i1, 0, i2);    // OK
int i3 = f("", s1, 0, i2);    // Ne compile pas
\end{lstlisting}
\end{onlyenv}
\end{frame}

\begin{frame}[fragile]
\frametitle{Utiliser le bon scope}
\begin{onlyenv}<+>
La même instance peut être utilisée avec différents types :
\begin{lstlisting}
class C {
 <T> void f(T t); // Utiliser un scope méthode
}
\end{lstlisting}
~\\
Exemple :
\begin{lstlisting}
class Reader {
 <T> T readFirst(List<T> list) {
  return list.get(0);
 }
}

Reader r = new Reader();
r.readFirst(Arrays.asList(1, 2, 3));       // 1
r.readFirst(Arrays.asList("a", "b", "c")); // a
\end{lstlisting}
\end{onlyenv}
\begin{onlyenv}<+>
La même instance doit être utilisée avec le même type :
\begin{lstlisting}
class C<T> { // Utiliser un scope classe
 void f(T t);
}
\end{lstlisting}
~\\
Exemple :
\begin{lstlisting}
class Reader<T> {
 T readFirst(List<T> list) {
  return list.get(0);
 }
}

Reader<Integer> r = new Reader<>();
r.readFirst(Arrays.asList(1, 2, 3));       // 1
r.readFirst(Arrays.asList("a", "b", "c")); // Ne compile pas
\end{lstlisting}
\end{onlyenv}
\begin{onlyenv}<+>
Si le type doit être partagé entre plusieurs méthodes :
\begin{itemize}
 \item Scope classe impose une instance par type
\begin{lstlisting}
class C<T> {
 void f1(T t1);
 void f2(T t2);
}
\end{lstlisting}
\end{itemize}
\end{onlyenv}
\end{frame}

\appendix
\newcounter{finalframe}
\setcounter{finalframe}{\value{framenumber}}

\scriptsize
\begin{thebibliography}{11}
\providecommand{\natexlab}[1]{#1}
\providecommand{\url}[1]{\texttt{#1}}
\expandafter\ifx\csname urlstyle\endcsname\relax
  \providecommand{\doi}[1]{doi: #1}\else
  \providecommand{\doi}{doi: \begingroup \urlstyle{rm}\Url}\fi

\bibitem[Ackerman(2011)]{ackerman_intelligence_2011}
P.~L. Ackerman.
\newblock Intelligence and {Expertise}.
\newblock In R.~J. Sternberg and S.~B. Kaufman, editors, \emph{The {Cambridge}
  {Handbook} of {Intelligence}}, pages 847--860. Cambridge University Press,
  Cambridge, 2011.
\newblock ISBN 978-0-511-97724-4.
\newblock URL
  \url{http://www.cambridge.org/us/academic/subjects/psychology/cognition/cambridge-handbook-intelligence}.

\bibitem[Davidson and Kemp(2011)]{davidson_contemporary_2011}
J.~E. Davidson and I.~A. Kemp.
\newblock Contemporary {Models} of {Intelligence}.
\newblock In R.~J. Sternberg and S.~B. Kaufman, editors, \emph{The {Cambridge}
  {Handbook} of {Intelligence}}, pages 58--82. Cambridge University Press,
  Cambridge, 2011.
\newblock ISBN 978-0-511-97724-4.
\newblock URL
  \url{http://www.cambridge.org/us/academic/subjects/psychology/cognition/cambridge-handbook-intelligence}.

\bibitem[Ericsson(2006)]{ericsson_introduction_2006}
K.~A. Ericsson.
\newblock An {Introduction} to {Cambridge} {Handbook} of {Expertise} and
  {Expert} {Performance}: {Its} {Development}, {Organization}, and {Content}.
\newblock In K.~A. Ericsson, N.~Charness, P.~J. Feltovich, and R.~R. Hoffman,
  editors, \emph{The {Cambridge} handbook of expertise and expert performance},
  pages 3--19. Cambridge University Press, New York, NY, US, 2006.
\newblock ISBN 0-521-60081-2 (Paperback); 0-521-84097-X (Hardcover);
  978-0-521-60081-1 (Paperback); 978-0-521-84097-2 (Hardcover).

\bibitem[Franklin(2014)]{franklin_history_2014}
S.~Franklin.
\newblock History, motivations, and core themes.
\newblock In K.~Frankish and W.~M. Ramsey, editors, \emph{The {Cambridge}
  handbook of artificial intelligence}, pages 15--33. Cambridge University
  Press, Cambridge, UK, 2014.
\newblock ISBN 978-0-521-87142-6.

\bibitem[Gottfredson(1997)]{gottfredson_mainstream_1997}
L.~S. Gottfredson.
\newblock Mainstream science on intelligence: {An} editorial with 52
  signatories, history, and bibliography.
\newblock \emph{Intelligence}, 24\penalty0 (1):\penalty0 13--23, Jan. 1997.
\newblock ISSN 01602896.
\newblock \doi{10.1016/S0160-2896(97)90011-8}.
\newblock URL
  \url{http://linkinghub.elsevier.com/retrieve/pii/S0160289697900118}.

\bibitem[Legg and Hutter(2007{\natexlab{a}})]{legg_collection_2007}
S.~Legg and M.~Hutter.
\newblock A {Collection} of {Definitions} of {Intelligence}.
\newblock In B.~Goertzel and P.~Wang, editors, \emph{Advances in artificial
  general intelligence: concepts, architectures and algorithms: proceedings of
  the {AGI} {Workshop} 2006}, number v. 157 in Frontiers in artificial
  intelligence and applications. IOS Press, Amsterdam ; Washington, DC,
  2007{\natexlab{a}}.
\newblock ISBN 978-1-58603-758-1.

\bibitem[Legg and Hutter(2007{\natexlab{b}})]{legg_universal_2007}
S.~Legg and M.~Hutter.
\newblock Universal {Intelligence}: {A} {Definition} of {Machine}
  {Intelligence}.
\newblock \emph{Minds and Machines}, 17\penalty0 (4):\penalty0 391--444, Nov.
  2007{\natexlab{b}}.
\newblock ISSN 0924-6495, 1572-8641.
\newblock \doi{10.1007/s11023-007-9079-x}.
\newblock URL \url{http://link.springer.com/10.1007/s11023-007-9079-x}.

\bibitem[Mackintosh(2011)]{mackintosh_history_2011}
N.~J. Mackintosh.
\newblock History of {Theories} and {Measurement} of {Intelligence}.
\newblock In R.~J. Sternberg and S.~B. Kaufman, editors, \emph{The {Cambridge}
  {Handbook} of {Intelligence}}, pages 3--19. Cambridge University Press,
  Cambridge, 2011.
\newblock ISBN 978-0-511-97724-4.
\newblock URL
  \url{http://www.cambridge.org/us/academic/subjects/psychology/cognition/cambridge-handbook-intelligence}.

\bibitem[Muehlhauser(2013)]{muehlhauser_what_2013}
L.~Muehlhauser.
\newblock What is {Intelligence}?, June 2013.
\newblock URL
  \url{https://intelligence.org/2013/06/19/what-is-intelligence-2/}.

\bibitem[Urbina(2011)]{urbina_tests_2011}
S.~Urbina.
\newblock Tests of {Intelligence}.
\newblock In R.~J. Sternberg and S.~B. Kaufman, editors, \emph{The {Cambridge}
  {Handbook} of {Intelligence}}, pages 20--38. Cambridge University Press,
  Cambridge, 2011.
\newblock ISBN 978-0-511-97724-4.
\newblock URL
  \url{http://www.cambridge.org/us/academic/subjects/psychology/cognition/cambridge-handbook-intelligence}.

\bibitem[Willis et~al.(2011)Willis, Dumont, and
  Kaufman]{willis_factor-analytic_2011}
J.~O. Willis, R.~Dumont, and A.~S. Kaufman.
\newblock Factor-{Analytic} {Models} of {Intelligence}.
\newblock In R.~J. Sternberg and S.~B. Kaufman, editors, \emph{The {Cambridge}
  {Handbook} of {Intelligence}}, pages 39--57. Cambridge University Press,
  Cambridge, 2011.
\newblock ISBN 978-0-511-97724-4.
\newblock URL
  \url{http://www.cambridge.org/us/academic/subjects/psychology/cognition/cambridge-handbook-intelligence}.

\end{thebibliography}

\normalsize

\begin{frame}
\frametitle{Summary}
\todo{todo}
\end{frame}

\setcounter{framenumber}{\value{finalframe}}
\end{document}
